% Luc Destombe --- licence CC BY-NC-SA 4.0
% https://creativecommons.org/licenses/by-nc-sa/4.0/deed.fr
% Fichier autonome : compiler avec pdflatex, deux fois (signets, nombre de pages).
\def\lyceeversion{1.0}
\documentclass[11pt,a4paper]{article}
\makeatletter
\newif\ifeleve
\elevefalse

\newif\iflycee@palatino
\lycee@palatinotrue

\RequirePackage[utf8]{inputenc}
\RequirePackage[T1]{fontenc}
\RequirePackage[french]{babel}
\RequirePackage{etoolbox}

\newcommand{\ShorthandsOff}{\@ifundefined{shorthandoff}{}{\shorthandoff{;:!?}}}
\newcommand{\ShorthandsOn}{\@ifundefined{shorthandon}{}{\shorthandon{;:!?}}}
\ShorthandsOff
\AtBeginDocument{\ShorthandsOn}
\AtBeginEnvironment{tikzpicture}{\ShorthandsOff}

\RequirePackage{geometry}
\RequirePackage{fancyhdr}
\RequirePackage{lastpage}
\RequirePackage{multicol}
\RequirePackage{array}
\RequirePackage{enumitem}
\RequirePackage{xcolor}
\RequirePackage{needspace}

\RequirePackage{amsmath,amssymb}
\RequirePackage{mathpazo}
\renewcommand{\ttdefault}{lmtt}
\DeclareMathSizes{10.95}{10.95}{8}{5.5}
\begingroup\catcode`\_=\active \catcode`\^=\active
\gdef_#1{\sb{\scriptscriptstyle #1}}
\gdef^#1{\sp{\scriptscriptstyle\lycee@moinsepais #1}}
\endgroup
\newcommand\lycee@moins{\mathbin{%
  \setbox\z@\hbox{$\m@th\scriptscriptstyle\lycee@moinsorig$}%
  \dimen@=.55\wd\z@ \advance\dimen@-.5pt
  \hbox to.75\wd\z@{\hss$\m@th\scriptscriptstyle\vcenter{\hbox{%
    \tikz\draw[line width=.5pt,line cap=round](0,0)--(\the\dimen@,0);}}$\hss}}}
\begingroup\lccode`\~=`\- \lowercase{\endgroup\let~}\lycee@moins
\newcommand\lycee@moinsepais{\mathcode`\-="8000 }
\AtBeginDocument{\mathchardef\lycee@moinsorig=\mathcode`\-\relax}
\AtBeginDocument{\catcode`\_=12 \mathcode`\_="8000
  \catcode`\^=12 \mathcode`\^="8000 }
\newcommand\lycee@exposants{%
  \fontdimen13\textfont2=.48\fontdimen6\textfont2
  \fontdimen14\textfont2=.48\fontdimen6\textfont2
  \fontdimen15\textfont2=.48\fontdimen6\textfont2
  \fontdimen13\scriptfont2=.48\fontdimen6\scriptfont2
  \fontdimen14\scriptfont2=.48\fontdimen6\scriptfont2
  \fontdimen15\scriptfont2=.48\fontdimen6\scriptfont2}
\every@math@size\expandafter{\the\every@math@size\lycee@exposants}
\newlength{\retraitcalculs}
\setlength{\retraitcalculs}{2em}

\RequirePackage{tikz}
\usetikzlibrary{arrows.meta,calc,babel,shapes.misc}
\RequirePackage[most]{tcolorbox}
\tcbuselibrary{skins,breakable}

\definecolor{coulDef}{HTML}{1F4E79}
\definecolor{coulProp}{HTML}{7030A0}
\definecolor{coulRem}{HTML}{C55A11}
\definecolor{coulEx}{HTML}{2E7D32}
\definecolor{coulExo}{HTML}{B03A2E}
\definecolor{coulPart}{HTML}{1F4E79}
\definecolor{coulMeth}{HTML}{1B6E4F}
\definecolor{coulCourbe}{HTML}{0B5ED7}
\definecolor{coulCourbeB}{HTML}{C00000}
\definecolor{coulCourbeC}{HTML}{2E7D32}
\definecolor{coulGrille}{HTML}{8C8C8C}
\definecolor{coulRep}{HTML}{1B6E4F}
\definecolor{coulBar}{HTML}{2E7D32}

\newcommand{\R}{\mathbb{R}}
\newcommand{\Rs}{\mathbb{R}^{*}}
\renewcommand{\le}{\leqslant}
\renewcommand{\ge}{\geqslant}
\renewcommand{\approx}{\simeq}
\newcommand{\vect}[1]{\overrightarrow{#1}}
\newcommand{\norme}[1]{\left\lVert #1 \right\rVert}
\newcommand{\intervalle}[2]{\left[#1\,;#2\right]}
\RequirePackage{cancel}
\renewcommand{\CancelColor}{\color{coulExo}}
\newcommand{\fois}[1]{\times{\color{coulExo}#1}}
\newcommand{\barre}[1]{\cancel{#1}}
\newlength{\lycee@eqh}\newlength{\lycee@eqd}
\newsavebox{\lycee@eqa}\newsavebox{\lycee@eqb}\newsavebox{\lycee@eqc}
\newcommand{\lycee@eqcol}[1]{\hbox{$\renewcommand{\arraystretch}{1.35}%
  \begin{array}[t]{@{}>{\displaystyle}r@{}>{\displaystyle{}}l@{}}#1\end{array}$}}
\newcommand{\lycee@eqmesure}[1]{%
  \ifdim\ht#1>\lycee@eqh\setlength{\lycee@eqh}{\ht#1}\fi
  \ifdim\dp#1>\lycee@eqd\setlength{\lycee@eqd}{\dp#1}\fi}
\newcommand{\lycee@eqtrait}{\hspace{1.2em}\tikz[baseline=0pt]\draw[dash pattern=on 2pt off 2pt,
  line width=0.5pt] (0,\lycee@eqh) -- (0,-\lycee@eqd);\hspace{1.2em}}
\newcommand{\eqproduit}[2]{\par\smallskip\noindent C'est une équation produit :
  $A\times B=0 \iff A=0$ ou $B=0$.\par\nopagebreak\smallskip
  \begingroup\setlength{\lycee@eqh}{0pt}\setlength{\lycee@eqd}{0pt}%
  \setbox\lycee@eqa=\lycee@eqcol{#1}\setbox\lycee@eqb=\lycee@eqcol{#2}%
  \lycee@eqmesure\lycee@eqa\lycee@eqmesure\lycee@eqb%
  \noindent\hspace*{\retraitcalculs}\usebox\lycee@eqa\lycee@eqtrait\usebox\lycee@eqb\endgroup\par\smallskip}
\newcommand{\eqproduittrois}[3]{\par\smallskip\noindent C'est une équation produit :
  $A\times B\times C=0 \iff A=0$ ou $B=0$ ou $C=0$.\par\nopagebreak\smallskip
  \begingroup\setlength{\lycee@eqh}{0pt}\setlength{\lycee@eqd}{0pt}%
  \setbox\lycee@eqa=\lycee@eqcol{#1}\setbox\lycee@eqb=\lycee@eqcol{#2}\setbox\lycee@eqc=\lycee@eqcol{#3}%
  \lycee@eqmesure\lycee@eqa\lycee@eqmesure\lycee@eqb\lycee@eqmesure\lycee@eqc%
  \noindent\hspace*{\retraitcalculs}\usebox\lycee@eqa\lycee@eqtrait\usebox\lycee@eqb\lycee@eqtrait\usebox\lycee@eqc\endgroup
  \par\smallskip}

\tikzset{grille/.style={coulGrille,line width=0.4pt}}
\newcommand{\quadrillage}[4]{\draw[grille,step=1] (#1,#2) grid (#3,#4);}
\tikzset{
  croix/.style={cross out,draw,line width=0.6pt,inner sep=0pt,minimum size=4pt},
  croixdroite/.style={croix,rotate=45,line width=0.8pt},
}
\newcommand{\pt}[3]{\path (#1) node[croix]{} node[#2,font=\small,inner sep=2.5pt]{$#3$};}

\newcommand{\lycee@repere}[4]{%
  \draw[grille,step=1] (#1,#2) grid (#3,#4);
  \draw[-{Stealth[length=2mm]},thick] (#1,0)--({#3+0.4},0) node[below left=-1pt]{$x$};
  \draw[-{Stealth[length=2mm]},thick] (0,#2)--(0,{#4+0.4}) node[below left=-1pt]{$y$};
  \node[below left,font=\scriptsize,inner sep=1.5pt] at (0,0) {$0$};}
\newcommand{\lycee@gradx}[1]{\foreach \k/\l in {#1}{%
  \draw (\k,0.07)--(\k,-0.07) node[below,font=\scriptsize,inner sep=1.5pt]{$\l$};}}
\newcommand{\lycee@grady}[1]{\foreach \k/\l in {#1}{%
  \draw (0.07,\k)--(-0.07,\k) node[left,font=\scriptsize,inner sep=1.5pt]{$\l$};}}
\newcommand{\lycee@intff}[2]{\left[#1\,;#2\right]}
\newcommand{\lycee@intfo}[2]{\left[#1\,;#2\right[}
\newcommand{\lycee@intof}[2]{\left]#1\,;#2\right]}
\newcommand{\lycee@intoo}[2]{\left]#1\,;#2\right[}
\newcommand{\lycee@ens}[1]{\left\{#1\right\}}
\AtBeginDocument{%
  \providecommand{\repere}{\lycee@repere}%
  \providecommand{\gradx}{\lycee@gradx}%
  \providecommand{\grady}{\lycee@grady}%
  \providecommand{\Cf}{\mathcal{C}\sb{\scriptscriptstyle f}}%
  \providecommand{\Cg}{\mathcal{C}\sb{\scriptscriptstyle g}}%
  \providecommand{\intff}{\lycee@intff}%
  \providecommand{\intfo}{\lycee@intfo}%
  \providecommand{\intof}{\lycee@intof}%
  \providecommand{\intoo}{\lycee@intoo}%
  \providecommand{\ens}{\lycee@ens}}

\newlist{questions}{enumerate}{3}
\setlist[questions,1]{label=\arabic*\textdegree),leftmargin=*,itemsep=2pt,topsep=3pt}
\setlist[questions,2]{label=\alph*),leftmargin=*,itemsep=1pt,topsep=2pt}
\setlist[questions,3]{label=\roman*),leftmargin=*,itemsep=1pt,topsep=1pt}
\setlist[itemize]{label=\textbullet,leftmargin=*,itemsep=2pt,topsep=3pt}

\newcommand{\besoinplace}[1]{\par\penalty\@M\begingroup
  \dimen@=#1\relax\dimen@ii=\pagegoal\advance\dimen@ii-\pagetotal
  \ifdim\dimen@>\dimen@ii\ifdim\dimen@ii>\z@\vfil\fi\break\fi\endgroup}

\newcommand{\lycee@pied}{\thepage/\pageref{LastPage}}
\newcommand{\entete}[2]{%
  \pagestyle{fancy}\fancyhf{}%
  \fancyhead[L]{\footnotesize\color{coulPart}\textbf{#1}}%
  \fancyhead[R]{\footnotesize\color{coulPart}\textbf{#2}}%
  \fancyfoot[C]{\footnotesize\lycee@pied}%
  \renewcommand{\headrulewidth}{0.4pt}%
  \renewcommand{\headrule}{\hbox to\headwidth{%
    \color{coulPart!40}\leaders\hrule height \headrulewidth\hfill}}}

\RequirePackage{ccicons}
\newcommand{\lycee@licence}{{\scriptsize\color{gray}%
  \copyright\ Luc Destombe --- \ccbyncsa\ CC BY-NC-SA 4.0}}
\newif\iflycee@licence \lycee@licencetrue
\newcommand{\sanslicence}{\lycee@licencefalse}
\AtBeginDocument{\iflycee@licence\AddToHookNext{shipout/foreground}{%
  \put(0,0){\rlap{\kern\dimexpr1in+\hoffset+\oddsidemargin+\textwidth\relax
    \llap{\raisebox{-\dimexpr1in+\voffset+\topmargin+\headheight+\headsep
      +\textheight+\footskip\relax}{\lycee@licence}}}}}\fi}

\newcommand{\lycee@site}{faire-des-maths.fr}
\newcommand{\lycee@version}{{\scriptsize\color{gray}\lycee@site
  \ifdefined\lyceeversion\ --- Version~\lyceeversion\fi}}
\AtBeginDocument{\iflycee@licence\AddToHookNext{shipout/foreground}{%
  \put(0,0){\rlap{\kern\dimexpr1in+\hoffset+\oddsidemargin\relax
    \raisebox{-\dimexpr1in+\voffset+\topmargin+\headheight+\headsep
      +\textheight+\footskip\relax}{\lycee@version}}}}\fi}

\newcommand{\formulesserrees}{%
  \setlength{\abovedisplayskip}{3pt}\setlength{\belowdisplayskip}{3pt}%
  \setlength{\abovedisplayshortskip}{2pt}\setlength{\belowdisplayshortskip}{3pt}}

\setlength{\parindent}{0pt}

\RequirePackage[implicit=false,hidelinks,pdfencoding=auto,psdextra,bookmarksopen,
  bookmarksopenlevel=1,pdfcreator={},pdfproducer={}]{hyperref}
\RequirePackage{bookmark}
\hypersetup{pdfauthor={Luc Destombe},pdfkeywords={CC BY-NC-SA 4.0}}
\newcounter{lycee@signet}
\newcommand{\signet}[3][]{\stepcounter{lycee@signet}%
  \edef\lycee@dest{\if\relax\detokenize{#1}\relax
    signet.\arabic{lycee@signet}\else#1\fi}%
  \hypertarget{\lycee@dest}{}\bookmark[level=#2,dest=\lycee@dest]{#3}}
\pdfstringdefDisableCommands{%
  \def\R{R}\def\vect#1{#1}\def\Cf{Cf}\def\Cg{Cg}\def\fois#1{×#1}%
  \def\barre#1{#1}\def\intervalle#1#2{[#1;#2]}\def\intff#1#2{[#1;#2]}%
  \def\textsuperscript#1{#1}\def\dfrac#1#2{#1/#2}\def\frac#1#2{#1/#2}%
  \def\vec#1{#1⃗}}%
\newcommand{\res}[1]{#1}
\geometry{top=2cm,bottom=1cm,left=1cm,right=1cm,headheight=14pt,footskip=0.6cm}

\RequirePackage{pgfplots}
\pgfplotsset{compat=1.18}
\RequirePackage{tkz-tab}

\newcounter{partiecours}
\renewcommand{\thepartiecours}{\Roman{partiecours}}
\newcommand{\partiecours}[1]{%
  \refstepcounter{partiecours}%
  \Needspace*{8\baselineskip}%
  \bigskip
  \noindent\begin{tcolorbox}[enhanced,colback=coulPart,colframe=coulPart,
      boxrule=0pt,arc=2pt,left=8pt,right=8pt,top=4pt,bottom=4pt,
      before skip=6pt,after skip=10pt]
    \color{white}\bfseries\large \thepartiecours{} --- #1%
    \signet[partie-\arabic{partiecours}]{1}{\thepartiecours{} --- #1}
  \end{tcolorbox}%
  \addcontentsline{toc}{section}{\thepartiecours{} --- #1}%
}

\newtcolorbox{boitecours}[3][]{%
  enhanced,breakable,
  colback=white,colframe=#2,coltitle=white,
  fonttitle=\bfseries,
  attach boxed title to top left={xshift=8pt,yshift=-\tcboxedtitleheight/2},
  boxed title style={colback=#2,arc=2pt,boxrule=0pt},
  boxrule=0.9pt,arc=2pt,
  left=8pt,right=8pt,top=8pt,bottom=6pt,
  title={#3},#1}

\newcounter{methode}
\newenvironment{definitionbox}[1][Définition]%
  {\begin{boitecours}[phantom={\signet{2}{#1}}]{coulDef}{#1}}{\end{boitecours}}
\newenvironment{proprietebox}[1][Propriété]%
  {\begin{boitecours}[phantom={\signet{2}{#1}}]{coulProp}{#1}}{\end{boitecours}}
\newenvironment{methodebox}[1][Méthode]{\stepcounter{methode}%
  \begin{boitecours}[phantom={\signet[methode-\arabic{methode}]{2}{#1}}]{coulMeth}{#1}}%
  {\end{boitecours}}

\newcommand{\etape}[1]{\par\smallskip\textbf{\color{coulMeth}#1}\par\nobreak\smallskip}
\newcommand{\enonce}[1]{\emph{#1}\par\smallskip}
\newcommand{\reponse}[1]{\par\smallskip
  {\footnotesize\itshape\color{black!55}Réponses : #1}}
\newcommand{\demo}[1]{\textbf{\color{coulMeth}Démonstration.} #1}
\newcommand{\afaire}[1]{%
  \par\smallskip
  \noindent{\small\color{coulExo}$\blacktriangleright$\ \textbf{À faire :}
    fiche d'exercices du chapitre, #1.}\par\smallskip}

\newenvironment{calculs}{\setlength{\jot}{5pt}\@fleqntrue\@mathmargin\retraitcalculs
  \setlength{\abovedisplayskip}{4pt}\setlength{\belowdisplayskip}{4pt}%
  \setlength{\abovedisplayshortskip}{4pt}\setlength{\belowdisplayshortskip}{4pt}%
  \csname alignat*\endcsname{2}}%
  {\csname endalignat*\endcsname}
\newcommand{\rem}[1]{\text{\small\itshape\color{black!60}#1}}
\newcommand{\explique}[2]{\par\smallskip\noindent
  \begin{minipage}[t]{0.57\linewidth}\raggedright #1\end{minipage}\hfill
  \begin{minipage}[t]{0.40\linewidth}\raggedright\leavevmode
    {\small\itshape\color{black!60}#2}\end{minipage}\par}
\newcommand{\cas}[1]{\par\smallskip\noindent\textbf{\color{coulExo}#1}\vadjust{\nobreak}}
\newcommand{\calc}[1]{\par\smallskip\textbf{\color{coulExo}#1}\par\nobreak}

\newtcolorbox{remarquebox}[1][Remarque]{%
  enhanced,breakable,colback=white,colframe=coulRem,
  boxrule=0pt,leftrule=2.5pt,arc=0pt,outer arc=0pt,
  left=8pt,right=6pt,top=5pt,bottom=5pt,
  before upper={\textbf{\color{coulRem}#1}\par\smallskip}}

\newtcolorbox{exemplebox}[1][Exemple]{%
  enhanced,breakable,colback=white,colframe=coulEx,
  boxrule=0pt,leftrule=2.5pt,arc=0pt,outer arc=0pt,
  left=8pt,right=6pt,top=5pt,bottom=5pt,
  before upper={\textbf{\color{coulEx}#1}\par\smallskip}}

\pgfplotsset{
  repere/.style={
    axis lines=middle,
    axis line style={-{Stealth[length=2.4mm]},thick},
    every axis x label/.style={at={(ticklabel* cs:1.0)},anchor=west},
    every axis y label/.style={at={(ticklabel* cs:1.0)},anchor=south},
    label style={font=\footnotesize},
    tick label style={font=\scriptsize},
    grid=both,
    major grid style={grille},
    minor tick num=0,
    tick align=inside,
    clip mode=individual,
  },
}
\tikzset{
  courbe/.style={coulCourbe,line width=1.1pt,smooth},
  courbeB/.style={coulCourbeB,line width=1.1pt,smooth},
  courbeC/.style={coulCourbeC,line width=1.1pt,smooth},
}

\newlist{etapes}{enumerate}{1}
\setlist[etapes]{label=\textbf{\arabic*.},leftmargin=*,itemsep=1pt,topsep=2pt}

\newlist{expressions}{enumerate}{1}
\setlist[expressions]{label={\Alph*\ $=$},leftmargin=*,itemsep=3pt,topsep=3pt}

\newcounter{exo}
\renewcommand{\theexo}{\ifnum\value{exo}<10 0\fi\arabic{exo}}
\newtcolorbox{exercice}[1][]{%
  enhanced,breakable,
  colback=white,colframe=coulExo,
  boxrule=0.9pt,arc=2pt,
  left=8pt,right=8pt,top=8pt,bottom=6pt,
  coltitle=white,fonttitle=\bfseries,
  attach boxed title to top left={xshift=8pt,yshift=-\tcboxedtitleheight/2},
  boxed title style={colback=coulExo,arc=2pt,boxrule=0pt},
  phantom={\signet{2}{Exercice \arabic{exo}}},
  title={Exercice \theexo\ifx\relax#1\relax\else{} --- #1\fi}}
\newenvironment{exo}[1][\relax]{\refstepcounter{exo}\begin{exercice}[#1]}{\end{exercice}}

  \newenvironment{correction}{\par}{\par}
  \newcommand{\autableau}{}
  \newcommand{\etapeexemples}{\etape{Exemples corrigés}}
  \newcommand{\etapetableau}[1]{\etape{#1}}
  \newenvironment{exempletableau}[1][Exemple]%
    {\begin{exemplebox}[#1]}{\end{exemplebox}}

\RequirePackage{comment}
\excludecomment{correctionavj}

\newcommand{\coursEntete}{}
\newcommand{\coursNiveau}{}
\newcommand{\coursTitre}{}
\newcommand{\coursSurTitre}{}
\newcommand{\coursSousTitre}{}

\newcommand{\enteteCours}[1]{\renewcommand{\coursEntete}{#1}}
\newcommand{\chapitrecours}[2]{\enteteCours{Chapitre #1 --- #2}}
\newcommand{\niveaucours}[1]{\renewcommand{\coursNiveau}{#1}}
\newcommand{\titrecours}[3][]{%
  \renewcommand{\coursTitre}{#2}%
  \renewcommand{\coursSurTitre}{#1}%
  \renewcommand{\coursSousTitre}{#3}}

\pagestyle{fancy}
\fancyhf{}
\fancyhead[L]{\footnotesize\color{coulPart}\textbf{\coursEntete}}
\fancyhead[R]{\footnotesize\color{coulPart}\textbf{\coursNiveau}}
\fancyfoot[C]{\footnotesize\thepage}
\renewcommand{\headrulewidth}{0.4pt}
\renewcommand{\headrule}{\hbox to\headwidth{\color{coulPart!40}\leaders\hrule height \headrulewidth\hfill}}

\setlength{\parskip}{2pt}

\newcommand{\lycee@mentioneleve}{\\[2pt]{\small\itshape Version élève}}
\AtBeginDocument{%
  \begin{center}
    \begin{tcolorbox}[enhanced,width=0.72\linewidth,colback=white,colframe=coulPart,
        boxrule=1.6pt,arc=3pt,halign=center,top=10pt,bottom=10pt]
      \ifx\coursSurTitre\empty
        {\Huge\bfseries\color{coulPart} \coursTitre}\\[4pt]
      \else
        {\Huge\bfseries\color{coulPart} \coursTitre}\\[3pt]
        {\Large\bfseries\color{coulPart} \coursSurTitre}\\[5pt]
      \fi
      {\normalsize \coursSousTitre}%
      
    \end{tcolorbox}
  \end{center}%
}

\RequirePackage{tikz}
\RequirePackage{xcolor}

\usetikzlibrary{decorations.text}
\definecolor{lycee@metal}{HTML}{1F4E79}
\definecolor{lycee@manche}{HTML}{8A5A2B}
\definecolor{lycee@etincelle}{HTML}{F28C28}
\newcommand{\lycee@pluriel}[2]{#1 #2\ifnum#1>1 s\fi}
\newcommand{\lycee@sceau}{%
  \edef\lycee@bas{\lycee@pluriel{\lyceerelectures}{relecture} et
    \lycee@pluriel{\lyceerodages}{rodage}}%
  \begin{tikzpicture}[line join=round]
    \fill[coulPart!10,draw=coulPart,line width=.8pt] (0:2.2)
      \foreach \i in {1,...,71}{ -- (\i*5:{isodd(\i)?2.1:2.2}) } -- cycle;
    \draw[coulPart,line width=1.2pt] (0,0) circle (1.93);
    \draw[coulPart,line width=.4pt] (0,0) circle (1.85);
    \draw[coulPart,line width=.4pt] (0,0) circle (1.2);
    \path[decorate,decoration={text along path,text align=center,
        text={|\fontsize{9.5}{11}\selectfont\bfseries\color{coulPart}|TEST{\'E} AVEC MES {\'E}L{\`E}VES}}]
      (180:1.45) arc (180:0:1.45);
    \path[decorate,decoration={text along path,text align=center,
        text/.expanded={|\noexpand\fontsize{8.5}{10}\noexpand\selectfont\noexpand\color{coulPart}|\lycee@bas}}]
      (195:1.68) arc (195:345:1.68);
    \begin{scope}[x=.25mm,y=-.25mm]
      \begin{scope}[shift={(-32,-40)}]%
        \fill[lycee@metal] (23,31) rectangle (41,49);
        \draw[lycee@etincelle,line width=.5pt,line cap=round]
          (20,30)--(15,27) (19,34)--(14,35) (44,30)--(49,27) (45,34)--(50,35);
        \fill[lycee@manche,rounded corners=.3mm] (38,24.6) rectangle (60,27.6);
        \fill[lycee@metal,rounded corners=.4mm] (25,20) rectangle (39,31);
        \fill[lycee@etincelle,rounded corners=.2mm] (25,20) rectangle (28,31);
      \end{scope}
    \end{scope}
    \ifdefined\lyceesceauplusieurs\def\lycee@annee{Ann{\'e}es~:}\else
      \def\lycee@annee{Ann{\'e}e}\fi
    \path[decorate,decoration={text along path,text align=center,
        text/.expanded={|\noexpand\fontsize{8}{10}\noexpand\selectfont\noexpand\color{coulPart}|%
          \unexpanded\expandafter{\lycee@annee} \unexpanded\expandafter{\lyceesceauannees}}}]
      (95:1.08) arc (95:445:1.08);
  \end{tikzpicture}}
\AtEndDocument{\ifdefined\lyceesceauannees
  \par\besoinplace{5cm}\vspace{1.5em}%
  \begin{center}\lycee@sceau\end{center}\fi}
\makeatother

\chapitrecours{10}{Fonctions de référence}
\niveaucours{Seconde}
\titrecours{FONCTIONS DE RÉFÉRENCE}{Cours --- Seconde}

\begin{document}

\newcommand{\schemacentre}[5]{%
  \begin{tikzpicture}[x=0.5cm,color=black,baseline=(base)]
    \coordinate (base) at (0,-0.1);
    \draw[-{Stealth[length=2mm]},thick] (-0.8,0)--(5.8,0);
    \ifnum#5=1 \draw[red,line width=1.6pt] (0,0)--(5,0);
    \else \path[red] (0,0) node[croix]{} (5,0) node[croix]{};\fi
    \draw (0,0.1)--(0,-0.1) node[below,font=\scriptsize\upshape]{$#1$};
    \draw (2.5,0.1)--(2.5,-0.1) node[below,font=\scriptsize\upshape]{$#2$};
    \draw (5,0.1)--(5,-0.1) node[below,font=\scriptsize\upshape]{$#3$};
    \draw[{Stealth[length=1.4mm]}-{Stealth[length=1.4mm]}] (0,0.3)--(2.5,0.3);
    \draw[{Stealth[length=1.4mm]}-{Stealth[length=1.4mm]}] (2.5,0.3)--(5,0.3);
    \node[font=\scriptsize\upshape,above] at (1.25,0.3) {$#4$};
    \node[font=\scriptsize\upshape,above] at (3.75,0.3) {$#4$};
  \end{tikzpicture}}

\partiecours{Fonction carré : courbe, variations, équations}

\begin{definitionbox}[Définition 1 --- Fonction carré]
La \textbf{fonction carré} est la fonction définie pour tout nombre réel $x$ par
${f(x)=x^{2}}$. Sa courbe s'appelle une \textbf{parabole}.

\smallskip
\begin{minipage}[c]{0.58\linewidth}\raggedright
Tableau de valeurs de $-3$ à $3$ avec un pas de $1$ (on ajoute $1$ à chaque fois) :
\begin{center}
\renewcommand{\arraystretch}{1.25}
\begin{tabular}{|c|*{7}{c|}}
\hline
$x$ & $-3$ & $-2$ & $-1$ & $0$ & $1$ & $2$ & $3$\\\hline
$x^{2}$ & $9$ & $4$ & $1$ & $0$ & $1$ & $4$ & $9$\\\hline
\end{tabular}
\end{center}
Deux nombres opposés ont le même carré : $(-3)^{2}=3^{2}=9$. La parabole est donc
\textbf{symétrique par rapport à l'axe des ordonnées}. Son point le plus bas est
l'origine, son \textbf{sommet}.
\end{minipage}\hfill
\begin{minipage}[c]{0.38\linewidth}\centering
\begin{tikzpicture}[x=0.36cm,y=0.36cm]
  \repere{-4}{-1}{4}{10}
  \gradx{-3,-2,-1,1,2,3} \grady{1,2,3,4,5,6,7,8,9}
  \draw[coulCourbe,line width=1.1pt,domain=-3.1:3.1,samples=60] plot (\x,{(\x)*(\x)});
  \foreach \a/\b in {-3/9,-2/4,-1/1,0/0,1/1,2/4,3/9} \path[coulCourbe] (\a,\b) node[croix]{};
  \node[coulCourbe,font=\footnotesize,right] at (2.4,8.6) {$y=x^{2}$};
\end{tikzpicture}
\end{minipage}
\end{definitionbox}

\begin{proprietebox}[Propriété 1 --- Variations et signe]
\begin{minipage}[c]{0.56\linewidth}\raggedright
La fonction carré est \textbf{décroissante} sur $\intof{-\infty}{0}$ et
\textbf{croissante} sur $\intfo{0}{+\infty}$. Son minimum est $0$, atteint en $0$ :
un carré est toujours positif ou nul.

\smallskip
Entre deux nombres \textbf{positifs}, l'ordre est \textbf{conservé} ; entre deux nombres
\textbf{négatifs}, il est \textbf{renversé}.
\end{minipage}\hfill
\begin{minipage}[c]{0.40\linewidth}\centering
\begin{tikzpicture}
  \tkzTabInit[lgt=1.2,espcl=1.6,deltacl=0.6]{$x$/1,$x^{2}$/1.6}{$-\infty$,$0$,$+\infty$}
  \tkzTabVar{+/ ,-/$0$,+/ }
\end{tikzpicture}
\end{minipage}

\smallskip
\textbf{Exemples d'abord.}
\explique{$2<3$ et $2^{2}<3^{2}$, car $4<9$.}{deux nombres positifs : même sens}
\explique{$-4<-3$, mais $(-4)^{2}>(-3)^{2}$, car $16>9$.}{deux nombres négatifs :
\textbf{attention}, changement de sens}

\smallskip
\demo{Sur $\intfo{0}{+\infty}$. On choisit deux nombres réels $a$ et $b$ tels que
$0<a<b$.}
\begin{calculs}
a-b &< 0 &\qquad& \rem{car $a<b$}\\
a+b &> 0 && \rem{somme de deux nombres positifs}\\
a^{2}-b^{2} &= (a+b)(a-b) && \rem{identité remarquable (chapitre 4)}\\
a^{2}-b^{2} &< 0 && \rem{produit d'un positif et d'un négatif}\\
a^{2} &< b^{2} &&
\end{calculs}
\explique{Conclusion : $0<a<b\Longrightarrow a^{2}<b^{2}$ ; la fonction carré est croissante
sur $\intfo{0}{+\infty}$.}{l'ordre est conservé}

\smallskip
\demo{Sur $\intof{-\infty}{0}$. On choisit deux nombres réels $a$ et $b$ tels que
$a<b<0$.}
\begin{calculs}
a-b &< 0 &\qquad& \rem{car $a<b$}\\
a+b &< 0 && \rem{somme de deux nombres négatifs}\\
a^{2}-b^{2} &= (a+b)(a-b) && \rem{identité remarquable}\\
a^{2}-b^{2} &> 0 && \rem{produit de deux nombres négatifs}\\
a^{2} &> b^{2} &&
\end{calculs}
\explique{Conclusion : $a<b<0\Longrightarrow a^{2}>b^{2}$ ; la fonction carré est décroissante
sur $\intof{-\infty}{0}$.}{l'ordre est renversé}
\end{proprietebox}

\begin{methodebox}[Méthode 1 --- Comparer deux carrés sans les calculer]
\etapeexemples
Comparer, sans les calculer :
\begin{questions}[label=\alph*)]
  \item $3{,}7^{2}$ et $4{,}2^{2}$
  \item $(-2{,}5)^{2}$ et $(-1{,}8)^{2}$
  \item $(-6)^{2}$ et $5{,}5^{2}$
\end{questions}
\begin{correction}
\cas{a)}
\explique{$0\le 3{,}7<4{,}2$ et la fonction carré est croissante sur
$\intfo{0}{+\infty}$, donc \res{$3{,}7^{2}<4{,}2^{2}$}.}{deux nombres positifs : l'ordre est
conservé}
\cas{b)}
\explique{$-2{,}5<-1{,}8\le 0$ et la fonction carré est décroissante sur
$\intof{-\infty}{0}$, donc \res{$(-2{,}5)^{2}>(-1{,}8)^{2}$}.}{deux nombres négatifs : l'ordre
est renversé}
\cas{c)}
\explique{Les deux nombres ne sont pas du même signe. Or $(-6)^{2}=6^{2}$.}{deux opposés ont
le même carré}
\explique{$0\le 5{,}5<6$, donc $5{,}5^{2}<6^{2}$, soit \res{$5{,}5^{2}<(-6)^{2}$}.}{on se
ramène à deux nombres positifs}
\end{correction}

\etape{À vous de jouer}
Comparer, sans les calculer :
\begin{center}
$1{,}3^{2}$ et $0{,}9^{2}$ \qquad $(-0{,}7)^{2}$ et $(-1{,}2)^{2}$ \qquad
$(-3{,}1)^{2}$ et $3{,}4^{2}$
\end{center}
\reponse{$1{,}3^{2}>0{,}9^{2}$ ;\; $(-0{,}7)^{2}<(-1{,}2)^{2}$ ;\; $(-3{,}1)^{2}<3{,}4^{2}$.}
\end{methodebox}

\afaire{exercices 1 à 3}

\begin{proprietebox}[Propriété 2 --- L'équation $x^{2}=k$ (rappel du chapitre 4)]
\begin{minipage}[c]{0.58\linewidth}\raggedright
\begin{itemize}
  \item Si $k>0$ : deux solutions, $-\sqrt{k}$ et $\sqrt{k}$.
  \item Si $k=0$ : une seule solution, $0$.
  \item Si $k<0$ : aucune solution, car un carré n'est jamais négatif.
\end{itemize}
Sur la parabole : la droite horizontale $y=k$ la coupe en deux points, en un point
(le sommet) ou pas du tout.
\end{minipage}\hfill
\begin{minipage}[c]{0.38\linewidth}\centering
\begin{tikzpicture}[x=0.36cm,y=0.36cm]
  \repere{-4}{-3}{4}{7}
  \gradx{-2,2} \grady{4}
  \draw[coulCourbe,line width=1.1pt,domain=-2.6:2.6,samples=60] plot (\x,{(\x)*(\x)});
  \draw[black!70,thick] (-4,4)--(4,4);
  \node[black!70,font=\scriptsize,above] at (3.2,4) {$y=4$};
  \draw[black!70,thick] (-4,-2)--(4,-2);
  \node[black!70,font=\scriptsize,below] at (3,-2) {$y=-2$};
  \path (-2,4) node[croix]{} (2,4) node[croix]{};
\end{tikzpicture}
\end{minipage}
\end{proprietebox}

\begin{methodebox}[Méthode 2 --- Résoudre une équation qui se ramène à $x^{2}=k$]
\etapeexemples
\begin{questions}[label=\alph*)]
  \item Résoudre $4x^{2}-5=x^{2}+22$.
  \item Résoudre $2x^{2}+7=5x^{2}+19$.
  \item Résoudre $6x^{2}-1=x^{2}+9$.
  \item Un jardin carré a une aire de $72$~m$^{2}$. Quelle est la longueur exacte de son
        côté ?
\end{questions}
\begin{correction}
\cas{a)}
\begin{calculs}
4x^{2}-5 &= x^{2}+22 &\qquad& \rem{on recopie l'équation}\\
4x^{2}-x^{2} &= 22+5 && \rem{les $x^{2}$ à gauche, les nombres à droite}\\
3x^{2} &= 27 &&\\
x^{2} &= 9 && \rem{on divise par $3$}
\end{calculs}
\explique{$9>0$, donc $S=\res{\ens{-3\,;3}}$.}{$\sqrt{9}=3$}
\cas{b)}
\begin{calculs}
2x^{2}+7 &= 5x^{2}+19 &\qquad& \rem{on recopie l'équation}\\
2x^{2}-5x^{2} &= 19-7 &&\\
-3x^{2} &= 12 &&\\
x^{2} &= -4 && \rem{on divise par $-3$}
\end{calculs}
\explique{Un carré n'est jamais négatif, donc $S=\res{\varnothing}$.}{$-4<0$}
\cas{c)}
\begin{calculs}
6x^{2}-1 &= x^{2}+9 &\qquad& \rem{on recopie l'équation}\\
5x^{2} &= 10 &&\\
x^{2} &= 2 &&
\end{calculs}
\explique{$2>0$, donc $S=\res{\ens{-\sqrt{2}\,;\sqrt{2}}}$.}{$2$ n'est pas un carré parfait :
on garde $\sqrt{2}$}
\cas{d)}
\explique{On note $x$ la longueur du côté, en mètres ($x>0$) : $x^{2}=72$.}{aire d'un carré :
côté $\times$ côté}
\explique{Deux solutions, $-\sqrt{72}$ et $\sqrt{72}$ ; une longueur est positive : on garde
$\sqrt{72}$.}{on revient au problème}
\begin{calculs}
\sqrt{72} &= \sqrt{36\times 2} &\qquad& \rem{$36$ est un carré parfait (chapitre 6)}\\
\sqrt{72} &= \sqrt{36}\times\sqrt{2} &&\\
\sqrt{72} &= 6\sqrt{2} &&
\end{calculs}
\explique{Le côté mesure \res{$6\sqrt{2}$~m}.}{valeur exacte}
\end{correction}

\etape{À vous de jouer}
\begin{questions}[label=\alph*)]
  \item Résoudre $5x^{2}-8=2x^{2}+40$.
  \item Résoudre $3x^{2}+10=x^{2}+2$.
  \item Résoudre $7x^{2}-4=3x^{2}+8$.
  \item Une dalle carrée a une aire de $50$~dm$^{2}$. Quelle est la longueur exacte de son
        côté ?
\end{questions}
\reponse{a) $S=\ens{-4\,;4}$ ;\; b)~$S=\varnothing$ ;\; c)~$S=\ens{-\sqrt3\,;\sqrt3}$ ;\;
 d)~$5\sqrt{2}$~dm.}
\end{methodebox}

\afaire{exercices 4 à 6}

\partiecours{Fonction carré : inéquations et résolution graphique}

\begin{proprietebox}[Propriété 3 --- Les inéquations $x^{2}\le k$ et $x^{2}\ge k$]
\begin{minipage}[c]{0.58\linewidth}\raggedright
Pour un nombre $k>0$ :
\begin{itemize}
  \item $x^{2}\le k$ pour les $x$ de $\intff{-\sqrt{k}}{\sqrt{k}}$ : la parabole est
        \textbf{sous} la droite $y=k$ (en rouge) ;
  \item $x^{2}\ge k$ pour les $x$ de $\intof{-\infty}{-\sqrt{k}}\cup\intfo{\sqrt{k}}{+\infty}$ :
        la parabole est \textbf{au-dessus} (en vert).
\end{itemize}
Avec $<$ ou $>$, les bornes $-\sqrt{k}$ et $\sqrt{k}$ sont exclues.
Si $k<0$ : $x^{2}\le k$ n'a aucune solution, et $x^{2}\ge k$ est vraie pour tout
nombre réel $x$.
\end{minipage}\hfill
\begin{minipage}[c]{0.38\linewidth}\centering
\begin{tikzpicture}[x=0.36cm,y=0.36cm]
  \repere{-4}{-1}{4}{9}
  \gradx{-2,2} \grady{4}
  \draw[coulCourbe,line width=1.1pt,domain=-2.9:2.9,samples=60] plot (\x,{(\x)*(\x)});
  \draw[black!70,thick] (-4,4)--(4,4);
  \node[black!70,font=\scriptsize,above] at (3.3,4) {$y=k$};
  \draw[black!70,dashed] (-2,4)--(-2,0) (2,4)--(2,0);
  \draw[red,line width=1.6pt] (-2,0)--(2,0);
  \draw[green!50!black,line width=1.6pt] (-4,0)--(-2,0) (2,0)--(4,0);
  \node[font=\scriptsize,below] at (-2.3,0) {$-\sqrt{k}$};
  \node[font=\scriptsize,below] at (2.2,0) {$\sqrt{k}$};
\end{tikzpicture}
\end{minipage}
\end{proprietebox}

\begin{methodebox}[Méthode 3 --- Résoudre une inéquation $x^{2}<k$, graphiquement puis exactement]
\etapeexemples
\begin{minipage}[c]{0.58\linewidth}\raggedright
La parabole $y=x^{2}$ et les droites $y=4$ et $y=9$ sont tracées ci-contre.
\begin{questions}[label=\alph*)]
  \item Résoudre graphiquement $x^{2}\le 4$, puis $x^{2}>9$.
  \item Résoudre $x^{2}<5$.
  \item Résoudre $3x^{2}-7\le 2x^{2}+9$.
\end{questions}
\end{minipage}\hfill
\begin{minipage}[c]{0.38\linewidth}\centering
\begin{tikzpicture}[x=0.36cm,y=0.36cm]
  \repere{-4}{-1}{4}{10}
  \gradx{-3,-2,-1,1,2,3} \grady{1,2,3,4,5,6,7,8,9}
  \draw[coulCourbe,line width=1.1pt,domain=-3.1:3.1,samples=60] plot (\x,{(\x)*(\x)});
  \draw[black!70,thick] (-4,4)--(4,4) (-4,9)--(4,9);
  \node[black!70,font=\scriptsize,above] at (-3.4,4) {$y=4$};
  \node[black!70,font=\scriptsize,above] at (-3.4,9) {$y=9$};
\end{tikzpicture}
\end{minipage}
\begin{correction}
\cas{a)}\par\nopagebreak
\begin{minipage}[c]{0.58\linewidth}\raggedright
La parabole est sur la droite $y=4$ ou au-dessous entre les abscisses $-2$ et $2$ :
${S=\res{\intff{-2}{2}}}$ (en rouge).

\smallskip
Elle est strictement au-dessus de la droite $y=9$ avant $-3$ et après $3$ :
${S=\res{\intoo{-\infty}{-3}\cup\intoo{3}{+\infty}}}$ (en vert).

\smallskip
{\small\itshape\color{black!60}On vérifie : $(-2)^{2}=2^{2}=4$ et $(-3)^{2}=3^{2}=9$ ;
« $>$ » : $-3$ et $3$ sont exclus.}
\end{minipage}\hfill
\begin{minipage}[c]{0.38\linewidth}\centering
\begin{tikzpicture}[x=0.36cm,y=0.36cm]
  \repere{-4}{-1}{4}{10}
  \gradx{-3,-2,-1,1,2,3} \grady{4,9}
  \draw[coulCourbe,line width=1.1pt,domain=-3.1:3.1,samples=60] plot (\x,{(\x)*(\x)});
  \draw[black!70,thick] (-4,4)--(4,4) (-4,9)--(4,9);
  \draw[black!70,dashed] (-2,4)--(-2,0) (2,4)--(2,0) (-3,9)--(-3,0) (3,9)--(3,0);
  \draw[red,line width=1.6pt] (-2,0)--(2,0);
  \draw[green!50!black,line width=1.6pt] (-4,0)--(-3,0) (3,0)--(4,0);
  \path (-2,4) node[croix]{} (2,4) node[croix]{} (-3,9) node[croix]{} (3,9) node[croix]{};
\end{tikzpicture}
\end{minipage}
\cas{b)}
\explique{$x^{2}=5$ a pour solutions $-\sqrt{5}$ et $\sqrt{5}$ : la droite $y=5$ coupe la
parabole en ces deux abscisses.}{$5$ n'est pas un carré parfait : valeurs exactes}
\explique{La parabole est strictement sous la droite $y=5$ entre ces deux abscisses :
${S=\res{\intoo{-\sqrt{5}}{\sqrt{5}}}}$.}{« $<$ » : bornes exclues}
\cas{c)}
\begin{calculs}
3x^{2}-7 &\le 2x^{2}+9 &\qquad& \rem{on recopie l'inéquation}\\
3x^{2}-2x^{2} &\le 9+7 && \rem{comme pour une équation}\\
x^{2} &\le 16 &&
\end{calculs}
\explique{$\sqrt{16}=4$, donc $S=\res{\intff{-4}{4}}$.}{« $\le$ » : bornes comprises}
\end{correction}

\etape{À vous de jouer}
\begin{questions}[label=\alph*)]
  \item À l'aide de la parabole ci-dessus, résoudre graphiquement $x^{2}<1$, puis
        $x^{2}\ge 4$.
  \item Résoudre $x^{2}\le 7$.
  \item Résoudre $2x^{2}+1>x^{2}+10$.
\end{questions}
\reponse{a) $S=\intoo{-1}{1}$ ;\; $S=\intof{-\infty}{-2}\cup\intfo{2}{+\infty}$ ;\;
 b)~$S=\intff{-\sqrt7}{\sqrt7}$ ;\; c)~$S=\intoo{-\infty}{-3}\cup\intoo{3}{+\infty}$.}
\end{methodebox}

\afaire{exercices 7 à 10}

\partiecours{Parabole et droite : position relative}

\begin{proprietebox}[Propriété 4 --- Position de deux courbes (rappel des chapitres 2 et 5)]
Pour savoir laquelle de deux courbes $y=f(x)$ et $y=g(x)$ est au-dessus de l'autre, on
étudie le \textbf{signe de la différence} $f(x)-g(x)$ :
\begin{itemize}
  \item là où $f(x)-g(x)\ge 0$, la courbe de $f$ est au-dessus de celle de $g$ ;
  \item là où $f(x)-g(x)\le 0$, elle est au-dessous ;
  \item là où $f(x)-g(x)=0$, les courbes se coupent.
\end{itemize}
Quand la différence est un produit, on fait son tableau de signes (chapitre 5).
\end{proprietebox}

\begin{methodebox}[Méthode 4 --- Étudier la position d'une parabole et d'une droite]
\etapeexemples
\begin{questions}[label=\alph*)]
  \item On considère la parabole $y=x^{2}$ et la droite (d) d'équation $y=x+2$.
        Développer $(x+1)(x-2)$, puis étudier la position de la parabole et de (d), et
        donner leurs points d'intersection.
  \item Pour $x\ge 0$, comparer $x^{2}$ et $x$ : étudier la position de la parabole
        $y=x^{2}$ et de la droite $y=x$ sur $\intfo{0}{+\infty}$.
\end{questions}
\begin{correction}
\cas{a)}
\begin{calculs}
(x+1)(x-2) &= x^{2}-2x+x-2 &\qquad& \rem{double distributivité (chapitre 4)}\\
(x+1)(x-2) &= x^{2}-x-2 &&
\end{calculs}
\explique{On étudie le signe de $x^{2}-(x+2)=x^{2}-x-2=(x+1)(x-2)$.}{la différence est le
produit qu'on vient de développer}
\explique{$x+1=0$ pour $x=-1$ et $x-2=0$ pour $x=2$ ; $a=1>0$ pour les deux : $-$ puis
$+$.}{signe de chaque facteur}
\begin{center}
\begin{tikzpicture}
  \tkzTabInit[lgt=3.4,espcl=2]{$x$/1,$x+1$/1,$x-2$/1,$x^{2}-(x+2)$/1}%
    {$-\infty$,$-1$,$2$,$+\infty$}
  \tkzTabLine{,-,z,+,t,+,}
  \tkzTabLine{,-,t,-,z,+,}
  \tkzTabLine{,+,z,-,z,+,}
\end{tikzpicture}
\end{center}
\begin{minipage}[c]{0.58\linewidth}\raggedright
La parabole est \res{au-dessous de (d) sur $\intff{-1}{2}$}, \res{au-dessus sur
$\intof{-\infty}{-1}$ et sur $\intfo{2}{+\infty}$}. {\small\itshape\color{black!60}(signes $-$, puis $+$)}

\smallskip

Elles se coupent pour $x=-1$ et $x=2$ ; $(-1)^{2}=1$ et $2^{2}=4$. {\small\itshape\color{black!60}(là où la
différence est nulle)}

\smallskip

Points d'intersection : \res{$(-1\,;1)$ et $(2\,;4)$}. {\small\itshape\color{black!60}(on vérifie sur (d) :
$-1+2=1$ et $2+2=4$)}

\smallskip

\end{minipage}\hfill
\begin{minipage}[c]{0.38\linewidth}\centering
\begin{tikzpicture}[x=0.36cm,y=0.36cm]
  \repere{-4}{-1}{4}{10}
  \gradx{-3,-2,-1,1,2,3} \grady{1,2,3,4,5,6,7,8,9}
  \draw[coulCourbe,line width=1.1pt,domain=-3.1:3.1,samples=60] plot (\x,{(\x)*(\x)});
  \draw[coulCourbeB,line width=1.1pt] (-3,-1)--(4,6);
  \node[coulCourbeB,font=\footnotesize,right] at (3.2,6.6) {(d)};
  \node[coulCourbe,font=\footnotesize,left] at (-2.7,8.6) {$y=x^{2}$};
  \path (-1,1) node[croix]{} (2,4) node[croix]{};
  \draw[red,line width=1.6pt] (-1,0)--(2,0);
\end{tikzpicture}
\end{minipage}
\cas{b)}
\explique{On étudie le signe de $x^{2}-x=x(x-1)$ sur $\intfo{0}{+\infty}$.}{on factorise par
$x$}
\begin{center}
\begin{tikzpicture}
  \tkzTabInit[lgt=3.4,espcl=2]{$x$/1,$x$/1,$x-1$/1,$x^{2}-x$/1}{$0$,$1$,$+\infty$}
  \tkzTabLine{z,+,t,+,}
  \tkzTabLine{,-,z,+,}
  \tkzTabLine{z,-,z,+,}
\end{tikzpicture}
\end{center}
\explique{Conclusion : \res{$x^{2}\le x$ sur $\intff{0}{1}$}, et \res{$x^{2}\ge x$ sur
$\intfo{1}{+\infty}$}.}{la parabole est sous la droite, puis au-dessus}
\end{correction}

\begin{remarquebox}[À retenir (démonstration du programme)]
Un nombre compris entre $0$ et $1$ est \textbf{plus grand que son carré} :
$0{,}5^{2}=0{,}25<0{,}5$. Au-delà de $1$, le carré est plus grand.
\end{remarquebox}

\etape{À vous de jouer}
\begin{questions}[label=\alph*)]
  \item On considère la parabole $y=x^{2}$ et la droite (d') d'équation $y=2x+3$.
        Développer $(x+1)(x-3)$, puis étudier la position de la parabole et de (d'), et
        donner leurs points d'intersection.
  \item Pour $x\ge 0$, étudier la position de la parabole $y=x^{2}$ et de la droite
        $y=3x$ sur $\intfo{0}{+\infty}$.
\end{questions}
\reponse{a) au-dessous de (d') sur $\intff{-1}{3}$, au-dessus ailleurs ;
$(-1\,;1)$ et $(3\,;9)$ ;\; b) au-dessous sur $\intff{0}{3}$, au-dessus sur
$\intfo{3}{+\infty}$.}
\end{methodebox}

\afaire{exercices 11 à 13}

\partiecours{Fonction cube}

\begin{definitionbox}[Définition 2 --- Fonction cube]
\begin{minipage}[c]{0.62\linewidth}\raggedright
La \textbf{fonction cube} est la fonction définie pour tout nombre réel $x$ par
${f(x)=x^{3}}$, avec $x^{3}=x\times x\times x$.

\smallskip
Tableau de valeurs de $-2$ à $2$ avec un pas de $1$ :
\begin{center}
\renewcommand{\arraystretch}{1.25}
\begin{tabular}{|c|*{5}{c|}}
\hline
$x$ & $-2$ & $-1$ & $0$ & $1$ & $2$\\\hline
$x^{3}$ & $-8$ & $-1$ & $0$ & $1$ & $8$\\\hline
\end{tabular}
\end{center}
Deux nombres opposés ont des cubes opposés : $(-2)^{3}=-8=-2^{3}$. La courbe est donc
\textbf{symétrique par rapport à l'origine} du repère.
\end{minipage}\hfill
\begin{minipage}[c]{0.34\linewidth}\centering
\begin{tikzpicture}[x=0.3cm,y=0.3cm]
  \repere{-3}{-9}{3}{9}
  \gradx{-2,-1,1,2} \grady{-8,-4,4,8}
  \draw[coulCourbe,line width=1.1pt,domain=-2.05:2.05,samples=60] plot (\x,{(\x)*(\x)*(\x)});
  \foreach \a/\b in {-2/-8,-1/-1,0/0,1/1,2/8} \path[coulCourbe] (\a,\b) node[croix]{};
  \node[coulCourbe,font=\footnotesize,left] at (1.8,8.4) {$y=x^{3}$};
\end{tikzpicture}
\end{minipage}
\end{definitionbox}

\begin{proprietebox}[Propriété 5 --- Variations, signe, équation $x^{3}=k$]
\begin{itemize}
  \item La fonction cube est \textbf{croissante} sur $\R$ : l'ordre est toujours conservé
        ($a<b$ donne $a^{3}<b^{3}$). Un cube a le signe du nombre : $(-2)^{3}<0$.
  \item Pour tout nombre réel $k$, l'équation $x^{3}=k$ a \textbf{une seule solution}, notée
        $\sqrt[3]{k}$ (racine cubique, touche de la calculatrice) : $x^{3}=8$ donne $x=2$,
        $x^{3}=5$ donne $x=\sqrt[3]{5}$.
  \item $x^{3}<k$ pour $x<\sqrt[3]{k}$, et $x^{3}>k$ pour $x>\sqrt[3]{k}$.
\end{itemize}
\end{proprietebox}

\begin{methodebox}[Méthode 5 --- Résoudre $x^{3}=k$ et $x^{3}<k$]
\etapeexemples
\begin{minipage}[c]{0.62\linewidth}\raggedright
\begin{questions}[label=\alph*)]
  \item Résoudre $x^{3}=64$, puis $x^{3}=-8$.
  \item Résoudre $2x^{3}+3=57$.
  \item À l'aide de la courbe ci-contre, résoudre graphiquement $x^{3}\le -1$.
  \item Résoudre $5x^{3}-4>3x^{3}+12$.
\end{questions}
\end{minipage}\hfill
\begin{minipage}[c]{0.34\linewidth}\centering
\begin{tikzpicture}[x=0.3cm,y=0.3cm]
  \repere{-3}{-9}{3}{9}
  \gradx{-2,-1,1,2} \grady{-8,-4,4,8}
  \draw[coulCourbe,line width=1.1pt,domain=-2.05:2.05,samples=60] plot (\x,{(\x)*(\x)*(\x)});
  \draw[black!70,thick] (-3,-1)--(3,-1);
  \node[black!70,font=\scriptsize,below] at (2.2,-1) {$y=-1$};
\end{tikzpicture}
\end{minipage}
\begin{correction}
\cas{a)}
\begin{calculs}
x^{3} &= 64 &\qquad& \rem{on recopie l'équation}\\
x &= \sqrt[3]{64} && \rem{une seule solution : la racine cubique}\\
x &= 4 && \rem{car $4^{3}=64$}
\end{calculs}
\explique{$S=\res{\ens{4}}$ ; de même, $x^{3}=-8$ donne $x=\sqrt[3]{-8}$, soit $x=-2$ :
$S=\res{\ens{-2}}$.}{$(-2)^{3}=-8$ : un cube peut être négatif}
\cas{b)}
\begin{calculs}
2x^{3}+3 &= 57 &\qquad& \rem{on recopie l'équation}\\
2x^{3} &= 54 &&\\
x^{3} &= 27 && \rem{on divise par $2$}\\
x &= \sqrt[3]{27} && \rem{racine cubique}\\
x &= 3 && \rem{car $3^{3}=27$}
\end{calculs}
\explique{$S=\res{\ens{3}}$.}{comme pour $x^{2}=k$, on isole d'abord le cube}
\cas{c)}\par\nopagebreak
\begin{minipage}[c]{0.62\linewidth}\raggedright
La courbe est sur la droite $y=-1$ ou au-dessous pour $x\le\sqrt[3]{-1}$, soit
$x\le -1$. {\small\itshape\color{black!60}($(-1)^{3}=-1$)}

\smallskip
La fonction cube est croissante, cela continue avant $-2$ :
${S=\res{\intof{-\infty}{-1}}}$. {\small\itshape\color{black!60}(la courbe ne remonte jamais)}
\end{minipage}\hfill
\begin{minipage}[c]{0.34\linewidth}\centering
\begin{tikzpicture}[x=0.3cm,y=0.3cm]
  \repere{-3}{-9}{3}{9}
  \gradx{-2,-1,1,2} \grady{-8,-4,4,8}
  \draw[coulCourbe,line width=1.1pt,domain=-2.05:2.05,samples=60] plot (\x,{(\x)*(\x)*(\x)});
  \draw[black!70,thick] (-3,-1)--(3,-1);
  \draw[black!70,dashed] (-1,-1)--(-1,0);
  \draw[red,line width=1.6pt] (-3,0)--(-1,0);
  \path (-1,-1) node[croix]{};
\end{tikzpicture}
\end{minipage}
\cas{d)}
\begin{calculs}
5x^{3}-4 &> 3x^{3}+12 &\qquad& \rem{on recopie l'inéquation}\\
2x^{3} &> 16 &&\\
x^{3} &> 8 &&\\
x &> \sqrt[3]{8} && \rem{la fonction cube conserve l'ordre}\\
x &> 2 && \rem{car $2^{3}=8$}
\end{calculs}
\explique{$S=\res{\intoo{2}{+\infty}}$.}{« $>$ » : $2$ est exclu}
\end{correction}

\etape{À vous de jouer}
\begin{questions}[label=\alph*)]
  \item Résoudre $x^{3}=1000$, puis $x^{3}=-1$.
  \item Résoudre $3x^{3}-5=19$.
  \item À l'aide de la courbe ci-dessus, résoudre graphiquement $x^{3}>8$.
  \item Résoudre $4x^{3}+1\le x^{3}-80$.
\end{questions}
\reponse{a) $S=\ens{10}$ ;\; $S=\ens{-1}$ ;\; b)~$S=\ens{2}$ ;\; c)~$S=\intoo{2}{+\infty}$ ;\;
 d)~$S=\intof{-\infty}{-3}$.}
\end{methodebox}

\afaire{exercices 14 à 17}

\partiecours{Fonction inverse}

\begin{definitionbox}[Définition 3 --- Fonction inverse]
La \textbf{fonction inverse} est la fonction définie pour tout nombre réel $x\ne 0$ par
${f(x)=\dfrac{1}{x}}$ : elle est définie sur $\intoo{-\infty}{0}\cup\intoo{0}{+\infty}$, car
$0$ n'a pas d'inverse. Sa courbe s'appelle une \textbf{hyperbole}.

\smallskip
\begin{minipage}[c]{0.60\linewidth}\raggedright
Tableau de valeurs de $-2$ à $2$ avec un pas de $0{,}5$, sans $0$ :
\begin{center}
\renewcommand{\arraystretch}{2.4}\setlength{\tabcolsep}{4pt}
\begin{tabular}{|c|*{8}{c|}}
\hline
$x$ & $-2$ & $-1{,}5$ & $-1$ & $-0{,}5$ & $0{,}5$ & $1$ & $1{,}5$ & $2$\\\hline
$\dfrac{1}{x}$ & $-\dfrac12$ & $-\dfrac23$ & $-1$ & $-2$ & $2$ & $1$ & $\dfrac23$
  & $\dfrac12$\\\hline
\end{tabular}
\end{center}
Deux nombres opposés ont des inverses opposés : l'hyperbole est
\textbf{symétrique par rapport à l'origine}.
\end{minipage}\hfill
\begin{minipage}[c]{0.36\linewidth}\centering
\begin{tikzpicture}[x=0.38cm,y=0.38cm]
  \repere{-4}{-4}{4}{4}
  \gradx{-2,-1,1,2} \grady{-2,-1,1,2}
  \draw[coulCourbe,line width=1.1pt,domain=0.26:4,samples=60] plot (\x,{1/(\x)});
  \draw[coulCourbe,line width=1.1pt,domain=-4:-0.26,samples=60] plot (\x,{1/(\x)});
  \foreach \a/\b in {-2/-0.5,-1/-1,-0.5/-2,0.5/2,1/1,2/0.5} \path[coulCourbe] (\a,\b) node[croix]{};
  \node[coulCourbe,font=\footnotesize,right] at (0.6,3.5) {$y=\frac{1}{x}$};
\end{tikzpicture}
\end{minipage}
\end{definitionbox}

\begin{proprietebox}[Propriété 6 --- Variations et signe]
\begin{minipage}[c]{0.56\linewidth}\raggedright
La fonction inverse est \textbf{décroissante} sur $\intoo{-\infty}{0}$ et
\textbf{décroissante} sur $\intoo{0}{+\infty}$ : sur chacun de ces intervalles, l'ordre
est renversé ($2<5$ donne $\frac12>\frac15$). Un inverse a le signe du nombre.
\end{minipage}\hfill
\begin{minipage}[c]{0.40\linewidth}\centering
\begin{tikzpicture}
  \tkzTabInit[lgt=1.2,espcl=1.6,deltacl=0.6]{$x$/1,$\frac{1}{x}$/1.6}{$-\infty$,$0$,$+\infty$}
  \tkzTabVar{+/ ,-D+/ / ,-/ }
\end{tikzpicture}
\end{minipage}

\smallskip
\demo{Exemple d'abord, avec $a=2$ et $b=5$ : $\frac12-\frac15=\frac{5}{10}-\frac{2}{10}
=\frac{3}{10}$, et $\frac{b-a}{ab}=\frac{3}{10}$. Soit maintenant deux nombres réels $a$ et
$b$ tels que $0<a<b$.}
\begin{calculs}
\frac{1}{a}-\frac{1}{b} &= \dfrac{1\fois{b}}{a\fois{b}}-\dfrac{1\fois{a}}{b\fois{a}}
  &\qquad& \rem{même dénominateur (chapitre 1)}\\
\frac{1}{a}-\frac{1}{b} &= \dfrac{b-a}{ab} &&\\
b-a &> 0 && \rem{car $a<b$}\\
ab &> 0 && \rem{produit de deux nombres positifs}
\end{calculs}
\explique{Donc $\frac{1}{a}-\frac{1}{b}>0$, soit $\frac{1}{a}>\frac{1}{b}$.}{quotient de
deux nombres positifs}
\explique{Conclusion : la fonction inverse est décroissante sur $\intoo{0}{+\infty}$ ; même
raisonnement sur $\intoo{-\infty}{0}$, où $ab>0$ aussi.}{l'ordre est renversé}
\end{proprietebox}

\begin{remarquebox}[Attention]
La fonction inverse n'est pas décroissante sur toute la réunion : $-1<1$ et
$\frac{1}{-1}<\frac{1}{1}$. On ne compare deux inverses que si les nombres ont le même
signe.
\end{remarquebox}

\begin{methodebox}[Méthode 6 --- Résoudre $\frac{1}{x}=k$ et $\frac{1}{x}\le k$]
\etapeexemples
\begin{questions}[label=\alph*)]
  \item Résoudre $\dfrac{1}{x}=4$, puis $\dfrac{1}{x}=-2$.
  \item Résoudre $\dfrac{1}{x}\le 2$ à l'aide d'un tableau de signes, puis vérifier sur
        l'hyperbole.
\end{questions}
\begin{correction}
\cas{a)}
\begin{calculs}
\frac{1}{x} &= 4 &\qquad& \rem{$x\ne 0$}\\
1 &= 4x && \rem{produit en croix}\\
x &= \frac14 &&
\end{calculs}
\explique{$S=\res{\ens{\frac14}}$ ; de même, $\frac{1}{x}=-2$ donne $1=-2x$, donc
$S=\res{\ens{-\frac12}}$.}{une seule solution : l'inverse de $k$}
\cas{b)}
\explique{On se ramène à une comparaison avec $0$ (chapitre 5).}{}
\begin{calculs}
\frac{1}{x}-2 &\le 0 &\qquad& \rem{on retire $2$}\\
\frac{1}{x}-\dfrac{2\fois{x}}{1\fois{x}} &\le 0 && \rem{même dénominateur}\\
\frac{1-2x}{x} &\le 0 &&
\end{calculs}
\explique{$1-2x=0$ pour $x=\frac12$, $a=-2<0$ : $+$ puis $-$. Valeur interdite
$0$.}{signe du numérateur ; double barre en $0$}
\begin{center}
\begin{tikzpicture}
  \tkzTabInit[lgt=2.6,espcl=2]{$x$/1,$1-2x$/1,$x$/1,$\frac{1-2x}{x}$/1.2}%
    {$-\infty$,$0$,$\frac12$,$+\infty$}
  \tkzTabLine{,+,t,+,z,-,}
  \tkzTabLine{,-,z,+,t,+,}
  \tkzTabLine{,-,d,+,z,-,}
\end{tikzpicture}
\end{center}
\begin{minipage}[c]{0.58\linewidth}\raggedright
$S=\res{\intoo{-\infty}{0}\cup\intfo{\frac12}{+\infty}}$. {\small\itshape\color{black!60}(signes $-$ et le $0$ ;
$0$ est exclu)}

\smallskip

L'hyperbole est bien sous la droite $y=2$ avant $0$ et à partir de
$\frac12$. {\small\itshape\color{black!60}(vérification graphique)}

\smallskip

\end{minipage}\hfill
\begin{minipage}[c]{0.38\linewidth}\centering
\begin{tikzpicture}[x=0.38cm,y=0.38cm]
  \repere{-4}{-4}{4}{4}
  \gradx{-2,-1,1,2} \grady{-2,-1,1,2}
  \draw[coulCourbe,line width=1.1pt,domain=0.26:4,samples=60] plot (\x,{1/(\x)});
  \draw[coulCourbe,line width=1.1pt,domain=-4:-0.26,samples=60] plot (\x,{1/(\x)});
  \draw[black!70,thick] (-4,2)--(4,2);
  \node[black!70,font=\scriptsize,above] at (3,2) {$y=2$};
  \draw[black!70,dashed] (0.5,2)--(0.5,0);
  \draw[red,line width=1.6pt] (-4,0)--(-0.08,0) (0.5,0)--(4,0);
  \path (0.5,2) node[croix]{};
\end{tikzpicture}
\end{minipage}
\end{correction}

\etape{À vous de jouer}
\begin{questions}[label=\alph*)]
  \item Résoudre $\dfrac{1}{x}=5$, puis $\dfrac{1}{x}=-4$.
  \item Résoudre $\dfrac{1}{x}\ge -1$ à l'aide d'un tableau de signes, puis vérifier sur
        l'hyperbole.
\end{questions}
\reponse{a) $S=\ens{\frac15}$ ;\; $S=\ens{-\frac14}$ ;\;
 b)~$S=\intof{-\infty}{-1}\cup\intoo{0}{+\infty}$.}
\end{methodebox}

\afaire{exercices 18 à 21}

\partiecours{Fonction racine carrée}

\begin{definitionbox}[Définition 4 --- Fonction racine carrée]
La \textbf{fonction racine carrée} est la fonction définie pour tout nombre réel $x\ge 0$
par ${f(x)=\sqrt{x}}$ : elle est définie sur $\intfo{0}{+\infty}$, car un nombre négatif n'a
pas de racine carrée (chapitre 6).

\smallskip
\begin{minipage}[c]{0.62\linewidth}\raggedright
Tableau de valeurs de $0$ à $9$ avec un pas de $1$ (valeurs arrondies au dixième) :
\begin{center}
\renewcommand{\arraystretch}{1.25}\setlength{\tabcolsep}{3.5pt}
\begin{tabular}{|c|*{10}{c|}}
\hline
$x$ & $0$ & $1$ & $2$ & $3$ & $4$ & $5$ & $6$ & $7$ & $8$ & $9$\\\hline
$\sqrt{x}$ & $0$ & $1$ & $1{,}4$ & $1{,}7$ & $2$ & $2{,}2$ & $2{,}4$ & $2{,}6$ & $2{,}8$
  & $3$\\\hline
\end{tabular}
\end{center}
La courbe est une demi-parabole « couchée » : c'est la symétrique de la moitié droite
de la parabole $y=x^{2}$ par rapport à la droite $y=x$, car $\sqrt{9}=3$ et $3^{2}=9$.
\end{minipage}\hfill
\begin{minipage}[c]{0.34\linewidth}\centering
\begin{tikzpicture}[x=0.38cm,y=0.38cm]
  \repere{-1}{-1}{10}{4}
  \gradx{1,4,9} \grady{1,2,3}
  \draw[coulCourbe,line width=1.1pt,domain=0:9.6,samples=80] plot (\x,{sqrt(\x)});
  \foreach \a/\b in {0/0,1/1,4/2,9/3} \path[coulCourbe] (\a,\b) node[croix]{};
  \node[coulCourbe,font=\footnotesize,above] at (6.5,2.6) {$y=\sqrt{x}$};
\end{tikzpicture}
\end{minipage}
\end{definitionbox}

\begin{proprietebox}[Propriété 7 --- Variations, équations et inéquations]
\begin{itemize}
  \item La fonction racine carrée est \textbf{croissante} sur $\intfo{0}{+\infty}$ : l'ordre
        est conservé. Une racine carrée est toujours positive ou nulle.
  \item $\sqrt{x}=k$ : si $k\ge 0$, une seule solution, $x=k^{2}$ ; si $k<0$, aucune.
  \item Pour $k\ge 0$ : $\sqrt{x}\le k$ pour les $x$ de $\intff{0}{k^{2}}$, et $\sqrt{x}\ge k$
        pour les $x$ de $\intfo{k^{2}}{+\infty}$.
\end{itemize}
Sans oublier : $x$ doit être positif, d'où la borne $0$.
\end{proprietebox}

\begin{methodebox}[Méthode 7 --- Résoudre $\sqrt{x}=k$ et $\sqrt{x}<k$]
\etapeexemples
\begin{minipage}[c]{0.60\linewidth}\raggedright
\begin{questions}[label=\alph*)]
  \item Résoudre $\sqrt{x}=6$, puis $\sqrt{x}=-3$.
  \item Résoudre $2\sqrt{x}-3=7$.
  \item Résoudre $\sqrt{x}\le 3$, et vérifier sur la courbe ci-contre.
  \item Résoudre $\sqrt{x}>1{,}5$.
\end{questions}
\end{minipage}\hfill
\begin{minipage}[c]{0.36\linewidth}\centering
\begin{tikzpicture}[x=0.38cm,y=0.38cm]
  \repere{-1}{-1}{10}{4}
  \gradx{1,4,9} \grady{1,2,3}
  \draw[coulCourbe,line width=1.1pt,domain=0:9.6,samples=80] plot (\x,{sqrt(\x)});
  \draw[black!70,thick] (-1,3)--(10,3);
  \node[black!70,font=\scriptsize,above] at (1.2,3) {$y=3$};
\end{tikzpicture}
\end{minipage}
\begin{correction}
\cas{a)}
\begin{calculs}
\sqrt{x} &= 6 &\qquad& \rem{$6\ge 0$}\\
(\sqrt{x})^{2} &= 6^{2} && \rem{on élève au carré}\\
x &= 36 && \rem{car $(\sqrt{x})^{2}=x$}
\end{calculs}
\explique{$S=\res{\ens{36}}$ ; une racine carrée n'est jamais négative, donc $\sqrt{x}=-3$
donne $S=\res{\varnothing}$.}{$-3<0$}
\cas{b)}
\begin{calculs}
2\sqrt{x}-3 &= 7 &\qquad& \rem{on recopie l'équation}\\
2\sqrt{x} &= 10 &&\\
\sqrt{x} &= 5 && \rem{on isole la racine}\\
(\sqrt{x})^{2} &= 5^{2} && \rem{on élève au carré}\\
x &= 25 &&
\end{calculs}
\explique{$S=\res{\ens{25}}$.}{on vérifie : $2\times 5-3=7$}
\cas{c)}\par\nopagebreak
\begin{minipage}[c]{0.60\linewidth}\raggedright
\begin{calculs}
\sqrt{x} &\le 3 &\qquad& \rem{deux nombres positifs}\\
(\sqrt{x})^{2} &\le 3^{2} && \rem{le carré conserve l'ordre}\\
x &\le 9 &&
\end{calculs}
Et $x\ge 0$, donc $S=\res{\intff{0}{9}}$. {\small\itshape\color{black!60}(on n'oublie pas la
borne $0$)}

\smallskip
La courbe est sous la droite $y=3$ de $0$ à $9$. {\small\itshape\color{black!60}(vérification
graphique)}
\end{minipage}\hfill
\begin{minipage}[c]{0.36\linewidth}\centering
\begin{tikzpicture}[x=0.38cm,y=0.38cm]
  \repere{-1}{-1}{10}{4}
  \gradx{1,4,9} \grady{1,2,3}
  \draw[coulCourbe,line width=1.1pt,domain=0:9.6,samples=80] plot (\x,{sqrt(\x)});
  \draw[black!70,thick] (-1,3)--(10,3);
  \draw[black!70,dashed] (9,3)--(9,0);
  \draw[red,line width=1.6pt] (0,0)--(9,0);
  \path (9,3) node[croix]{};
\end{tikzpicture}
\end{minipage}
\cas{d)}
\begin{calculs}
\sqrt{x} &> 1{,}5 &\qquad& \rem{on recopie l'inéquation}\\
(\sqrt{x})^{2} &> 1{,}5^{2} && \rem{le carré conserve l'ordre}\\
x &> 2{,}25 &&
\end{calculs}
\explique{$S=\res{\intoo{2{,}25}{+\infty}}$.}{« $>$ » : borne exclue}
\end{correction}

\etape{À vous de jouer}
\begin{questions}[label=\alph*)]
  \item Résoudre $\sqrt{x}=7$, puis $\sqrt{x}=-1$.
  \item Résoudre $3\sqrt{x}+2=14$.
  \item Résoudre $\sqrt{x}<4$.
  \item Résoudre $\sqrt{x}\ge 2{,}5$.
\end{questions}
\reponse{a) $S=\ens{49}$ ;\; $S=\varnothing$ ;\; b)~$S=\ens{16}$ ;\; c)~$S=\intfo{0}{16}$ ;\;
 d)~$S=\intfo{6{,}25}{+\infty}$.}
\end{methodebox}

\afaire{exercices 22 à 25}

\partiecours{Valeur absolue d'un nombre}

\begin{definitionbox}[Définition 5 --- Valeur absolue, distance]
\begin{itemize}
  \item La \textbf{valeur absolue} d'un nombre réel $a$, notée $|a|$, est la
        \textbf{distance} entre $a$ et $0$ sur la droite numérique : $\lvert -4|=4$ et $|3|=3$.
        Si $a\ge 0$, $|a|=a$ ; si $a<0$, $|a|=-a$ (on retire le signe $-$).
  \item La \textbf{distance} entre deux nombres réels $a$ et $b$ est $|a-b|$, qui vaut aussi
        $|b-a|$ : la distance entre $-2$ et $5$ vaut $|5-(-2)|=7$.
\end{itemize}
\begin{center}
\begin{tikzpicture}[x=0.6cm]
  \draw[-{Stealth[length=2.4mm]},thick] (-3.5,0)--(6.5,0);
  \foreach \k in {-3,...,6} \draw (\k,0.1)--(\k,-0.1) node[below,font=\scriptsize]{$\k$};
  \draw[coulCourbe,line width=1.4pt,{Stealth[length=2mm]}-{Stealth[length=2mm]}] (-2,0.35)--(5,0.35);
  \node[coulCourbe,font=\footnotesize,above] at (1.5,0.35) {distance $7$};
  \path (-2,0) node[croix]{} (5,0) node[croix]{};
\end{tikzpicture}
\end{center}
\end{definitionbox}

\begin{proprietebox}[Propriété 8 --- Racine carrée d'un carré]
Pour tout nombre réel $a$ : $\sqrt{a^{2}}=|a|$. Par exemple,
$\sqrt{(-3)^{2}}=\sqrt{9}$, soit $3$, c'est-à-dire $\lvert -3|$ : ce n'est pas $-3$, car une racine carrée est positive.
\end{proprietebox}

\begin{methodebox}[Méthode 8 --- Calculer une valeur absolue, une distance]
\etapeexemples
\begin{questions}[label=\alph*)]
  \item Calculer $\lvert -7|$ et $|2{,}5|$.
  \item Calculer la distance entre $-3$ et $4$.
  \item Écrire sans valeur absolue $|3-\sqrt{10}|$.
  \item Écrire sans racine carrée $\sqrt{(2-\sqrt{5})^{2}}$.
\end{questions}
\begin{correction}
\cas{a)}
\explique{$\lvert -7|=\res{7}$ et $|2{,}5|=\res{2{,}5}$.}{distance à $0$}
\cas{b)}
\begin{calculs}
|4-(-3)| &= |4+3| &\qquad& \rem{distance : $|b-a|$}\\
|4-(-3)| &= \res{7} &&
\end{calculs}
\cas{c)}
\explique{$\sqrt{10}>\sqrt{9}=3$, donc $3-\sqrt{10}<0$.}{on cherche le signe}
\explique{$|3-\sqrt{10}|=\res{\sqrt{10}-3}$.}{nombre négatif : on prend son opposé}
\cas{d)}
\begin{calculs}
\sqrt{(2-\sqrt{5})^{2}} &= |2-\sqrt{5}| &\qquad& \rem{propriété 8}
\end{calculs}
\explique{$\sqrt{5}>\sqrt{4}=2$, donc $2-\sqrt{5}<0$ et le résultat vaut
\res{$\sqrt{5}-2$}.}{opposé de $2-\sqrt{5}$}
\end{correction}

\etape{À vous de jouer}
\begin{questions}[label=\alph*)]
  \item Calculer $\lvert -12|$ et $|0{,}4|$.
  \item Calculer la distance entre $-6$ et $-1$.
  \item Écrire sans valeur absolue $|\sqrt{2}-1|$.
  \item Écrire sans racine carrée $\sqrt{(3-\sqrt{11})^{2}}$.
\end{questions}
\reponse{a) $12$ ;\; $0{,}4$ ;\; b)~$5$ ;\; c)~$\sqrt2-1$ ;\; d)~$\sqrt{11}-3$.}
\end{methodebox}

\afaire{exercices 26 à 29}

\partiecours{Intervalle de centre $a$ et de rayon $r$}

\begin{proprietebox}[Propriété 9 --- Intervalle de centre $a$ et de rayon $r$]
\begin{minipage}[c]{0.56\linewidth}\raggedright
Pour $r>0$ : $|x-a|\le r$ veut dire que $x$ est à une distance au plus $r$ de $a$,
c'est-à-dire $a-r\le x\le a+r$. Les solutions forment l'intervalle
$\intff{a-r}{a+r}$, de \textbf{centre} $a$ et de \textbf{rayon} $r$.
\end{minipage}\hfill
\begin{minipage}[c]{0.40\linewidth}\centering
\begin{tikzpicture}[x=0.55cm]
  \draw[-{Stealth[length=2.4mm]},thick] (-0.5,0)--(6.5,0);
  \draw[red,line width=1.6pt] (1,0)--(5,0);
  \foreach \k/\l in {1/a-r,3/a,5/a+r} \draw (\k,0.1)--(\k,-0.1) node[below,font=\scriptsize]{$\l$};
  \draw[font=\scriptsize] (1,0.3)--(1,0.3) (1.1,0.35) node[above right]{$r$}
    (3.1,0.35) node[above right]{$r$};
  \draw[{Stealth[length=1.6mm]}-{Stealth[length=1.6mm]}] (1,0.3)--(3,0.3);
  \draw[{Stealth[length=1.6mm]}-{Stealth[length=1.6mm]}] (3,0.3)--(5,0.3);
\end{tikzpicture}
\end{minipage}
\end{proprietebox}

\begin{methodebox}[Méthode 9 --- Résoudre $|x-a|\le r$ ; traduire un intervalle]
\etapeexemples
\begin{questions}[label=\alph*)]
  \item Résoudre $|x-3|\le 2$, puis $|x+1|\le 4$.
  \item Résoudre $|x-2|=5$.
  \item Écrire l'intervalle $\intff{4}{10}$ à l'aide d'une valeur absolue.
  \item Une pièce doit mesurer $50$~mm, à $0{,}2$~mm près. Écrire cette condition sur sa
        longueur $L$ avec une valeur absolue, puis donner l'intervalle des longueurs
        acceptées.
\end{questions}
\begin{correction}
\cas{a)}
\explique{Le centre vaut $3$ et le rayon $2$ ; $3-2=1$ et $3+2=5$, donc
$S=\res{\intff{1}{5}}$.}{les $x$ à distance au plus $2$ de $3$\par
\schemacentre{1}{3}{5}{2}{1}}
\explique{$|x+1|=|x-(-1)|$ : centre $-1$, rayon $4$, donc
$S=\res{\intff{-5}{3}}$.}{$x+1=x-(-1)$\par\schemacentre{-5}{-1}{3}{4}{1}}
\cas{b)}
\explique{Les nombres à distance exactement $5$ de $2$ sont $2-5$ et $2+5$ :
$S=\res{\ens{-3\,;7}}$.}{les deux bornes seulement\par\schemacentre{-3}{2}{7}{5}{0}}
\cas{c)}
\begin{calculs}
a &= \frac{4+10}{2} &\qquad& \rem{le centre est le milieu}\\
a &= 7 &&\\
r &= 10-7 && \rem{distance du centre à une borne}\\
r &= 3 &&
\end{calculs}
\explique{Donc $x\in\intff{4}{10}$ s'écrit \res{$|x-7|\le 3$}.}{\schemacentre{4}{7}{10}{3}{1}}
\cas{d)}
\explique{La distance entre $L$ et $50$ vaut au plus $0{,}2$ : \res{$|L-50|\le 0{,}2$}.}{on
traduit « à $0{,}2$ près »}
\explique{$L\in\res{\intff{49{,}8}{50{,}2}}$.}{$50-0{,}2$ et $50+0{,}2$\par
\schemacentre{49{,}8}{50}{50{,}2}{0{,}2}{1}}
\end{correction}

\etape{À vous de jouer}
\begin{questions}[label=\alph*)]
  \item Résoudre $|x-6|\le 1$, puis $|x+3|\le 2$.
  \item Résoudre $|x-1|=4$.
  \item Écrire l'intervalle $\intff{-2}{8}$ à l'aide d'une valeur absolue.
  \item Un sachet de bonbons doit peser $250$~g, à $5$~g près. Écrire cette condition sur sa
        masse $m$ avec une valeur absolue, puis donner l'intervalle des masses acceptées.
\end{questions}
\reponse{a) $S=\intff{5}{7}$ ;\; $S=\intff{-5}{-1}$ ;\; b)~$S=\ens{-3\,;5}$ ;\;
 c)~$|x-3|\le 5$ ;\; d)~$|m-250|\le 5$, $\intff{245}{255}$.}
\end{methodebox}

\afaire{exercices 30 à 32}

\partiecours{Fonction valeur absolue}

\begin{definitionbox}[Définition 6 --- Fonction valeur absolue]
\begin{minipage}[c]{0.60\linewidth}\raggedright
La \textbf{fonction valeur absolue} est la fonction définie pour tout nombre réel $x$ par
${f(x)=|x|}$ : $|x|=x$ si $x\ge 0$, et $|x|=-x$ si $x<0$.

\smallskip
Tableau de valeurs de $-3$ à $3$ avec un pas de $1$ :
\begin{center}
\renewcommand{\arraystretch}{1.25}
\begin{tabular}{|c|*{7}{c|}}
\hline
$x$ & $-3$ & $-2$ & $-1$ & $0$ & $1$ & $2$ & $3$\\\hline
$|x|$ & $3$ & $2$ & $1$ & $0$ & $1$ & $2$ & $3$\\\hline
\end{tabular}
\end{center}
Sa courbe est formée de deux demi-droites (chapitre 9) : $y=-x$ pour $x\le 0$, et $y=x$
pour $x\ge 0$. Elle est symétrique par rapport à l'axe des ordonnées, comme la parabole.
\end{minipage}\hfill
\begin{minipage}[c]{0.36\linewidth}\centering
\begin{tikzpicture}[x=0.38cm,y=0.38cm]
  \repere{-4}{-1}{4}{4}
  \gradx{-3,-2,-1,1,2,3} \grady{1,2,3}
  \draw[coulCourbe,line width=1.1pt] (-3.6,3.6)--(0,0)--(3.6,3.6);
  \foreach \a/\b in {-3/3,-2/2,-1/1,0/0,1/1,2/2,3/3} \path[coulCourbe] (\a,\b) node[croix]{};
  \node[coulCourbe,font=\footnotesize,right] at (1.4,3.4) {$y=|x|$};
\end{tikzpicture}
\end{minipage}
\end{definitionbox}

\begin{proprietebox}[Propriété 10 --- Variations, signe, équations et inéquations]
\begin{minipage}[c]{0.56\linewidth}\raggedright
\begin{itemize}
  \item Décroissante sur $\intof{-\infty}{0}$, croissante sur $\intfo{0}{+\infty}$ ;
        minimum $0$ en $0$, donc $|x|\ge 0$.
  \item Pour $k>0$ : $|x|=k$ a deux solutions, $-k$ et $k$ ; $|x|\le k$ pour les $x$ de
        $\intff{-k}{k}$ ; $|x|\ge k$ pour les $x$ de
        $\intof{-\infty}{-k}\cup\intfo{k}{+\infty}$.
\end{itemize}
\end{minipage}\hfill
\begin{minipage}[c]{0.40\linewidth}\centering
\begin{tikzpicture}
  \tkzTabInit[lgt=1.2,espcl=1.6,deltacl=0.6]{$x$/1,$|x|$/1.6}{$-\infty$,$0$,$+\infty$}
  \tkzTabVar{+/ ,-/$0$,+/ }
\end{tikzpicture}
\end{minipage}

\smallskip
Lien avec la fonction carré : $\sqrt{x^{2}}=|x|$, donc $x^{2}\le 9$ et $|x|\le 3$ ont les
mêmes solutions, $\intff{-3}{3}$.
\end{proprietebox}

\begin{methodebox}[Méthode 10 --- Résoudre $|x|=k$ et $|x|<k$]
\etapeexemples
\begin{minipage}[c]{0.60\linewidth}\raggedright
\begin{questions}[label=\alph*)]
  \item Résoudre $|x|=2{,}5$, puis $|x|=-2$.
  \item À l'aide de la courbe ci-contre, résoudre graphiquement $|x|<3$, puis $|x|\ge 1$.
  \item Résoudre $2|x|-1\le 5$.
\end{questions}
\end{minipage}\hfill
\begin{minipage}[c]{0.36\linewidth}\centering
\begin{tikzpicture}[x=0.38cm,y=0.38cm]
  \repere{-4}{-1}{4}{4}
  \gradx{-3,-2,-1,1,2,3} \grady{1,2,3}
  \draw[coulCourbe,line width=1.1pt] (-3.6,3.6)--(0,0)--(3.6,3.6);
  \draw[black!70,thick] (-4,3)--(4,3) (-4,1)--(4,1);
  \node[black!70,font=\scriptsize,above] at (-3.3,3) {$y=3$};
  \node[black!70,font=\scriptsize,above] at (-3.3,1) {$y=1$};
\end{tikzpicture}
\end{minipage}
\begin{correction}
\cas{a)}
\explique{Les nombres à distance $2{,}5$ de $0$ sont solutions, donc $S=\res{\ens{-2{,}5\,;2{,}5}}$.}{deux
solutions opposées}
\explique{Une valeur absolue n'est jamais négative : $S=\res{\varnothing}$.}{une distance est
positive}
\cas{b)}\par\nopagebreak
\begin{minipage}[c]{0.60\linewidth}\raggedright
La courbe est strictement sous la droite $y=3$ entre $-3$ et $3$ :
${S=\res{\intoo{-3}{3}}}$ (en rouge). {\small\itshape\color{black!60}(« $<$ » : bornes exclues)}

\smallskip

Elle est sur la droite $y=1$ ou au-dessus avant $-1$ et après $1$ :
${S=\res{\intof{-\infty}{-1}\cup\intfo{1}{+\infty}}}$ (en vert). {\small\itshape\color{black!60}(« $\ge$ » : bornes
comprises)}

\smallskip

\end{minipage}\hfill
\begin{minipage}[c]{0.36\linewidth}\centering
\begin{tikzpicture}[x=0.38cm,y=0.38cm]
  \repere{-4}{-1}{4}{4}
  \gradx{-3,-2,-1,1,2,3} \grady{1,3}
  \draw[coulCourbe,line width=1.1pt] (-3.6,3.6)--(0,0)--(3.6,3.6);
  \draw[black!70,thick] (-4,3)--(4,3) (-4,1)--(4,1);
  \draw[black!70,dashed] (-3,3)--(-3,0) (3,3)--(3,0);
  \draw[red,line width=1.6pt] (-3,0)--(3,0);
  \draw[green!50!black,line width=1.6pt] (-4,-0.25)--(-1,-0.25) (1,-0.25)--(4,-0.25);
  \path (-3,3) node[croix]{} (3,3) node[croix]{} (-1,1) node[croix]{} (1,1) node[croix]{};
\end{tikzpicture}
\end{minipage}
\cas{c)}
\begin{calculs}
2|x|-1 &\le 5 &\qquad& \rem{on recopie l'inéquation}\\
2|x| &\le 6 &&\\
|x| &\le 3 && \rem{on isole la valeur absolue}
\end{calculs}
\explique{Donc $S=\res{\intff{-3}{3}}$.}{les nombres à distance au plus $3$ de $0$}
\end{correction}

\etape{À vous de jouer}
\begin{questions}[label=\alph*)]
  \item Résoudre $|x|=4$, puis $|x|=-1$.
  \item À l'aide de la courbe ci-dessus, résoudre graphiquement $|x|\le 1$, puis $|x|>3$.
  \item Résoudre $3|x|+2=14$.
\end{questions}
\reponse{a) $S=\ens{-4\,;4}$ ;\; $S=\varnothing$ ;\; b)~$S=\intff{-1}{1}$ ;\;
$S=\intoo{-\infty}{-3}\cup\intoo{3}{+\infty}$ ;\; c)~$S=\ens{-4\,;4}$.}
\end{methodebox}

\afaire{exercices 33 à 35}

\begin{remarquebox}[Synthèse --- Les courbes à connaître]
\newcommand{\minicourbe}[3]{%
  \begin{minipage}[t]{0.32\linewidth}\centering
  \begin{tikzpicture}[x=0.32cm,y=0.32cm]
    \repere{-3}{-3}{3}{3}
    #1
  \end{tikzpicture}\par\footnotesize #2\par #3
  \end{minipage}}
\minicourbe{\draw[coulCourbe,line width=1.1pt] (-2.5,-3)--(3,2.5);}{affine $y=ax+b$}{droite ;
croissante si $a>0$}\hfill
\minicourbe{\draw[coulCourbe,line width=1.1pt,domain=-1.7:1.7,samples=40] plot (\x,{(\x)*(\x)});}%
{carré $y=x^{2}$}{définie sur $\R$ ; décroissante, puis croissante}\hfill
\minicourbe{\draw[coulCourbe,line width=1.1pt,domain=-1.42:1.42,samples=40] plot (\x,{(\x)*(\x)*(\x)});}%
{cube $y=x^{3}$}{définie sur $\R$ ; croissante}

\medskip
\minicourbe{\draw[coulCourbe,line width=1.1pt,domain=0.34:3,samples=40] plot (\x,{1/(\x)});
  \draw[coulCourbe,line width=1.1pt,domain=-3:-0.34,samples=40] plot (\x,{1/(\x)});}%
{inverse $y=\frac{1}{x}$}{définie pour $x\ne 0$ ; décroissante sur chaque intervalle}\hfill
\minicourbe{\draw[coulCourbe,line width=1.1pt,domain=0:3,samples=40] plot (\x,{sqrt(\x)});}%
{racine carrée $y=\sqrt{x}$}{définie pour $x\ge 0$ ; croissante}\hfill
\minicourbe{\draw[coulCourbe,line width=1.1pt] (-2.8,2.8)--(0,0)--(2.8,2.8);}%
{valeur absolue $y=|x|$}{définie sur $\R$ ; décroissante, puis croissante}
\end{remarquebox}

\end{document}
